\chapter[Trabalhos Relacionados]{Trabalhos Relacionados}

A adoção de tecnologias baseadas em blockchain em cenários que envolvem dados sensíveis, como prontuários médicos, levanta desafios importantes relacionados à privacidade, conformidade regulatória e arquitetura de sistemas.

Nesse contexto, \cite{GDPRBlockchainTension} apresenta uma revisão sistemática da literatura que analisa a compatibilidade entre blockchains públicas e a General Data Protection Regulation (GDPR). O estudo identifica três problemas estruturais centrais: (i) a imutabilidade dos dados, que dificulta a implementação do direito ao esquecimento; (ii) a transparência e rastreabilidade, que podem expor dados pessoais sensíveis; e (iii) a descentralização, que dificulta a definição de responsabilidades legais. Esses aspectos evidenciam tensões fundamentais entre os princípios da blockchain e os requisitos regulatórios de proteção de dados.

No domínio específico de sistemas de saúde, trabalhos como \cite{BlockchainVehiclesForEfficientMedicalRecordManagement} e \cite{MedicalRecordSystemUsingBlockchainBigDataAndTokenization} destacam limitações dos modelos tradicionais, incluindo a centralização dos dados, baixa interoperabilidade entre sistemas, vulnerabilidades de segurança e ineficiência no compartilhamento de informações. Como alternativa, ambos propõem o uso de blockchain para garantir integridade, rastreabilidade e descentralização dos registros médicos, aumentando a confiabilidade e a transparência no gerenciamento dos dados.

Entretanto, embora a utilização de blockchain mitigue diversos problemas estruturais, ela também introduz desafios relacionados à conformidade com legislações de proteção de dados, como a Lei Geral de Proteção de Dados (LGPD). Nesse sentido, \cite{BlockGuardLGPDHyperledgerFabric} propõe uma abordagem estruturada para adequação de soluções baseadas em Hyperledger Fabric aos requisitos legais brasileiros. O trabalho apresenta um \textit{framework} baseado no mapeamento entre os princípios da LGPD e as funcionalidades da plataforma, organizado em cinco domínios principais: gestão do ciclo de vida dos dados, controle de consentimento, operacionalização dos direitos dos titulares, mecanismos de segurança e governança, e documentação de conformidade.

Entre as contribuições do modelo proposto, destacam-se o uso de uma arquitetura híbrida, combinando armazenamento \textit{on-chain} e \textit{off-chain}, a implementação de \textit{smart contracts} para controle granular de consentimento e auditoria, e a definição de mecanismos para viabilizar o direito ao esquecimento por meio de técnicas como eliminação lógica e criptográfica. Além disso, o framework propõe um processo metodológico em fases que orienta a implementação contínua da conformidade regulatória.

Dessa forma, os trabalhos analisados demonstram que, embora a blockchain apresente vantagens significativas em termos de integridade, transparência e descentralização, sua adoção em sistemas que manipulam dados pessoais exige adaptações arquiteturais e mecanismos adicionais para garantir conformidade com legislações de proteção de dados.
